\section{Límite de una función}

\begin{definition}[Límite]
  Sean \(A \subset \R^n\), \(x_0 \in A'\) (punto de acumulación de \(A\)) y \(f\colon A \to \R^k\) una función. Decimos que \(y \in \R^k\) es el \emph{límite} de \(f\) en \(x_0\), y escribimos \(\lim_{x \to x_0} f(x) = y\), si
  \[
    (\forall \e > 0)(\exists \delta > 0) \colon (\forall x \in A) \quad 0 < \norm{x - x_0} < \delta \Rightarrow \norm{f(x) - y} < \e.
  \]
\end{definition}

\begin{lemma}
  Si \(A \subset \R^n\) y \(x_0 \in A'\), una función \(f\colon A \to \R^k\) tiene a lo sumo un límite en \(x_0\).
\end{lemma}

\begin{proof}
  Sean \(y, y' \in \R^k\) límites de \(f\) en \(x_0\) y supongamos \(y \neq y'\). Sea \(\e = \frac{1}{2} \norm{y - y'} > 0\). Existen \(\delta_1, \delta_2 > 0\) tales que
  \[
    0 < \norm{x - x_0} < \delta_1 \Rightarrow \norm{f(x) - y} < \e, \qquad 0 < \norm{x - x_0} < \delta_2 \Rightarrow \norm{f(x) - y'} < \e.
  \]
  Sea \(\delta = \min\{\delta_1, \delta_2\}\). Como \(x_0 \in A'\), existe \(x \in A\) con \(0 < \norm{x - x_0} < \delta\), y para ese \(x\),
  \[
    \norm{y - y'} \leq \norm{y - f(x)} + \norm{f(x) - y'} < 2\e = \norm{y - y'},
  \]
  lo cual es absurdo. Luego \(y = y'\).
\end{proof}

\begin{prop}
  Sean \(A \subset \R^n\), \(x_0 \in A'\), \(f\colon A \to \R^k\) e \(y \in \R^k\). Son equivalentes:
  \begin{enumerate}[label=(\arabic*)]
    \item \(\lim_{x \to x_0} f(x) = y\).
    \item Para todo \(\e > 0\) existe \(\delta > 0\) tal que \(f\bigl( B_{\delta}(x_0) \cap (A \setminus \{x_0\}) \bigr) \subset B_{\e}(y)\).
    \item Para todo entorno \(V\) de \(y\) existe un entorno \(U\) de \(x_0\) tal que \(f\bigl( U \cap (A \setminus \{x_0\}) \bigr) \subset V\).
  \end{enumerate}
\end{prop}

\begin{proof}
  (1) \(\Rightarrow\) (2) Dado \(\e > 0\), sea \(\delta > 0\) como en la definición. Si \(x \in B_{\delta}(x_0) \cap (A \setminus \{x_0\})\), entonces \(0 < \norm{x - x_0} < \delta\), luego \(\norm{f(x) - y} < \e\), es decir, \(f(x) \in B_{\e}(y)\).

  (2) \(\Rightarrow\) (3) Sea \(V\) un entorno de \(y\). Existe \(\e > 0\) con \(B_{\e}(y) \subset V\). Por (2) existe \(\delta > 0\) con \(f\bigl( B_{\delta}(x_0) \cap (A \setminus \{x_0\}) \bigr) \subset B_{\e}(y) \subset V\). Basta tomar \(U = B_{\delta}(x_0)\), que es un entorno de \(x_0\).

  (3) \(\Rightarrow\) (1) Sea \(\e > 0\). Como \(B_{\e}(y)\) es un entorno de \(y\), existe un entorno \(U\) de \(x_0\) con \(f\bigl( U \cap (A \setminus \{x_0\}) \bigr) \subset B_{\e}(y)\). Como \(U\) es entorno de \(x_0\), existe \(\delta > 0\) con \(B_{\delta}(x_0) \subset U\). Si \(x \in A\) cumple \(0 < \norm{x - x_0} < \delta\), entonces \(x \in U \cap (A \setminus \{x_0\})\), luego \(f(x) \in B_{\e}(y)\), es decir, \(\norm{f(x) - y} < \e\).
\end{proof}

\begin{prop}[Caracterización por sucesiones]
  Sean \(A \subset \R^n\), \(x_0 \in A'\), \(f\colon A \to \R^k\) e \(y \in \R^k\). Entonces \(\lim_{x \to x_0} f(x) = y\) si y sólo si para toda sucesión \({(x_m)}_{m \in \N} \subset A \setminus \{x_0\}\) con \(x_m \to x_0\) se cumple \(f(x_m) \to y\).
\end{prop}

\begin{proof}
  (\(\Rightarrow\)) Sea \({(x_m)} \subset A \setminus \{x_0\}\) con \(x_m \to x_0\) y sea \(\e > 0\). Existe \(\delta > 0\) tal que \(0 < \norm{x - x_0} < \delta \Rightarrow \norm{f(x) - y} < \e\). Como \(x_m \to x_0\), existe \(m_0\) con \(\norm{x_m - x_0} < \delta\) para \(m > m_0\), y como \(x_m \neq x_0\), resulta \(\norm{f(x_m) - y} < \e\) para \(m > m_0\).

  (\(\Leftarrow\)) Por contrarrecíproco, supongamos que \(\lim_{x \to x_0} f(x) \neq y\). Entonces existe \(\e > 0\) tal que para todo \(\delta > 0\) hay un \(x \in A\) con \(0 < \norm{x - x_0} < \delta\) y \(\norm{f(x) - y} \geq \e\). Tomando \(\delta = \frac{1}{m}\) obtenemos \(x_m \in A \setminus \{x_0\}\) con \(\norm{x_m - x_0} < \frac{1}{m}\), de modo que \(x_m \to x_0\), pero \(\norm{f(x_m) - y} \geq \e\) para todo \(m\), es decir, \(f(x_m) \not\to y\).
\end{proof}

\begin{prop}[Álgebra de límites]
  Sean \(A \subset \R^n\), \(x_0 \in A'\), \(f, g\colon A \to \R^k\) y \(h\colon A \to \R\).
  \begin{enumerate}[label=(\arabic*)]
    \item Si \(\lim_{x \to x_0} f(x) = y\) y \(\lim_{x \to x_0} g(x) = z\), entonces \(\lim_{x \to x_0} \bigl( f(x) + g(x) \bigr) = y + z\).
    \item Si \(\lim_{x \to x_0} h(x) = r\) y \(\lim_{x \to x_0} f(x) = y\), entonces \(\lim_{x \to x_0} h(x) f(x) = r y\).
    \item Si \(\lim_{x \to x_0} h(x) = r \neq 0\) y \(h(x) \neq 0\) para todo \(x \in A\), entonces \(\lim_{x \to x_0} \dfrac{1}{h(x)} = \dfrac{1}{r}\).
  \end{enumerate}
\end{prop}

\begin{proof}
  Sea \({(x_m)} \subset A \setminus \{x_0\}\) con \(x_m \to x_0\). Por la caracterización por sucesiones, \(f(x_m) \to y\), \(g(x_m) \to z\) y \(h(x_m) \to r\). Por las propiedades de límites de sucesiones, \(f(x_m) + g(x_m) \to y + z\), \(h(x_m) f(x_m) \to r y\) y \(\frac{1}{h(x_m)} \to \frac{1}{r}\). Como esto vale para toda tal sucesión, se obtienen (1), (2) y (3).
\end{proof}

\begin{prop}
  Sean \(A \subset \R^n\), \(x_0 \in A'\), \(f\colon A \to \R^k\) y \(\lim_{x \to x_0} f(x) = y\).
  \begin{enumerate}[label=(\arabic*)]
    \item Si \(z \neq y\), existe \(r > 0\) tal que \(z \notin f\bigl( B_r(x_0) \cap (A \setminus \{x_0\}) \bigr)\).
    \item Si \(k = 1\) e \(y > 0\), existe \(r > 0\) tal que \(f(x) > 0\) para todo \(x \in B_r(x_0) \cap (A \setminus \{x_0\})\).
    \item Existe \(r > 0\) tal que \(f\) es acotada en \(B_r(x_0) \cap A\), es decir, el conjunto \(\{ \norm{f(x)} : x \in B_r(x_0) \cap A \}\) está acotado superiormente.
  \end{enumerate}
\end{prop}

\begin{proof}
  (1) Como \(\norm{y - z} > 0\), tomando \(\e = \norm{y - z}\) existe \(r > 0\) tal que \(0 < \norm{x - x_0} < r \Rightarrow \norm{f(x) - y} < \norm{y - z}\). En particular \(f(x) \neq z\) para tales \(x\).

  (2) Tomando \(\e = y > 0\), existe \(r > 0\) con \(|f(x) - y| < y\), es decir, \(f(x) > 0\), para todo \(x \in A\) con \(0 < \norm{x - x_0} < r\).

  (3) Tomando \(\e = 1\), existe \(r > 0\) tal que para \(x \in A\) con \(0 < \norm{x - x_0} < r\) se tiene \(\norm{f(x)} \leq \norm{f(x) - y} + \norm{y} < 1 + \norm{y}\). Si \(x_0 \notin A\), entonces \(1 + \norm{y}\) es cota de \(\norm{f(x)}\) en \(B_r(x_0) \cap A\). Si \(x_0 \in A\), el punto \(x_0\) no está contemplado por la definición de límite, y una cota es \(\max\{ 1 + \norm{y}, \norm{f(x_0)} \}\).
\end{proof}

\begin{theorem}[Teorema del sándwich]
  Sean \(A \subset \R^n\), \(x_0 \in A'\) y \(f, g, h\colon A \to \R\) tales que \(f(x) \leq g(x) \leq h(x)\) para todo \(x \in A\). Si \(\lim_{x \to x_0} f(x) = \lim_{x \to x_0} h(x) = y\), entonces \(\lim_{x \to x_0} g(x) = y\).
\end{theorem}

\begin{proof}
  Dado \(\e > 0\), existen \(\delta_1, \delta_2 > 0\) tales que
  \[
    0 < \norm{x - x_0} < \delta_1 \Rightarrow |f(x) - y| < \e, \qquad 0 < \norm{x - x_0} < \delta_2 \Rightarrow |h(x) - y| < \e.
  \]
  Sea \(\delta = \min\{\delta_1, \delta_2\}\). Si \(x \in A\) cumple \(0 < \norm{x - x_0} < \delta\), entonces
  \[
    y - \e < f(x) \leq g(x) \leq h(x) < y + \e,
  \]
  es decir, \(|g(x) - y| < \e\). Luego \(\lim_{x \to x_0} g(x) = y\).
\end{proof}
